\documentclass[10pt,a4paper]{amsart}
\usepackage{amsmath,amssymb,amsthm,graphicx}
% For commutative diagrams we use Xy-pic
\usepackage[all]{xy}

% We use 2cell for 2-commutative diagrams.
\xyoption{2cell}
% Theorem environments.
%
\theoremstyle{plain}
\newtheorem{theorem}[subsection]{Theorem}
\newtheorem{proposition}[subsection]{Proposition}
\newtheorem{lemma}[subsection]{Lemma}

\theoremstyle{definition}
\newtheorem{definition}[subsection]{Definition}
\newtheorem{example}[subsection]{Example}
\newtheorem{exercise}[subsection]{Exercise}
\newtheorem{situation}[subsection]{Situation}

\theoremstyle{remark}
\newtheorem{remark}[subsection]{Remark}
\newtheorem{remarks}[subsection]{Remarks}

\numberwithin{equation}{subsection}

% Macros
%
\def\lim{\mathop{\mathrm{lim}}\nolimits}
\def\colim{\mathop{\mathrm{colim}}\nolimits}
\def\Spec{\mathop{\mathrm{Spec}}}
\def\Hom{\mathop{\mathrm{Hom}}\nolimits}
\def\Ext{\mathop{\mathrm{Ext}}\nolimits}
\def\SheafHom{\mathop{\mathcal{H}\!\mathit{om}}\nolimits}
\def\SheafExt{\mathop{\mathcal{E}\!\mathit{xt}}\nolimits}
\def\Sch{\mathit{Sch}}
\def\Mor{\mathop{\mathrm{Mor}}\nolimits}
\def\Ob{\mathop{\mathrm{Ob}}\nolimits}
\def\Sh{\mathop{\mathit{Sh}}\nolimits}
\def\NL{\mathop{N\!L}\nolimits}
\def\CH{\mathop{\mathrm{CH}}\nolimits}
\def\proetale{{pro\text{-}\acute{e}tale}}
\def\etale{{\acute{e}tale}}
\def\QCoh{\mathit{QCoh}}
\def\Ker{\mathop{\mathrm{Ker}}}
\def\Im{\mathop{\mathrm{Im}}}
\def\Coker{\mathop{\mathrm{Coker}}}
\def\Coim{\mathop{\mathrm{Coim}}}

% Boxtimes
%
\DeclareMathSymbol{\boxtimes}{\mathbin}{AMSa}{"02}

%
% Macros for moduli stacks/spaces
%
\def\QCohstack{\mathcal{QC}\!\mathit{oh}}
\def\Cohstack{\mathcal{C}\!\mathit{oh}}
\def\Spacesstack{\mathcal{S}\!\mathit{paces}}
\def\Quotfunctor{\mathrm{Quot}}
\def\Hilbfunctor{\mathrm{Hilb}}
\def\Curvesstack{\mathcal{C}\!\mathit{urves}}
\def\Polarizedstack{\mathcal{P}\!\mathit{olarized}}
\def\Complexesstack{\mathcal{C}\!\mathit{omplexes}}
% \Pic is the operator that assigns to X its picard group, usage \Pic(X)
% \Picardstack_{X/B} denotes the Picard stack of X over B
% \Picardfunctor_{X/B} denotes the Picard functor of X over B
\def\Pic{\mathop{\mathrm{Pic}}\nolimits}
\def\Picardstack{\mathcal{P}\!\mathit{ic}}
\def\Picardfunctor{\mathrm{Pic}}
\def\Deformationcategory{\mathcal{D}\!\mathit{ef}}
\newenvironment{reference}{}{}
\newenvironment{slogan}{}{}
\newenvironment{history}{}{}
\usepackage{cleveref}
\crefname{originalsection}{original section}{original sections}
\Crefname{originalsection}{Original section}{Original sections}
\makeatletter
\newlabel{resolve-section-minus-one}{{0C2I}{source}{}{sourcepointer-1}{}}
\newlabel{resolve-section-minus-one@cref}{{[originalsection][0][]0C2I}{[1][0][]source}{}{}{}}
\makeatother
\crefname{originallemma}{original lemma}{original lemmas}
\Crefname{originallemma}{Original lemma}{Original lemmas}
\makeatletter
\newlabel{curves-lemma-nodal}{{0C4D}{source}{}{sourcepointer-2}{}}
\newlabel{curves-lemma-nodal@cref}{{[originallemma][0][]0C4D}{[1][0][]source}{}{}{}}
\makeatother
\crefname{originallemma}{original lemma}{original lemmas}
\Crefname{originallemma}{Original lemma}{Original lemmas}
\makeatletter
\newlabel{curves-lemma-smooth-relative-dimension-1}{{0CD7}{source}{}{sourcepointer-3}{}}
\newlabel{curves-lemma-smooth-relative-dimension-1@cref}{{[originallemma][0][]0CD7}{[1][0][]source}{}{}{}}
\makeatother
\crefname{originallemma}{original lemma}{original lemmas}
\Crefname{originallemma}{Original lemma}{Original lemmas}
\makeatletter
\newlabel{morphisms-lemma-smooth-at-point}{{01V9}{source}{}{sourcepointer-4}{}}
\newlabel{morphisms-lemma-smooth-at-point@cref}{{[originallemma][0][]01V9}{[1][0][]source}{}{}{}}
\makeatother
\crefname{originallemma}{original lemma}{original lemmas}
\Crefname{originallemma}{Original lemma}{Original lemmas}
\makeatletter
\newlabel{algebra-lemma-regular-ring-CM}{{00NQ}{source}{}{sourcepointer-5}{}}
\newlabel{algebra-lemma-regular-ring-CM@cref}{{[originallemma][0][]00NQ}{[1][0][]source}{}{}{}}
\makeatother
\crefname{originallemma}{original lemma}{original lemmas}
\Crefname{originallemma}{Original lemma}{Original lemmas}
\makeatletter
\newlabel{algebra-lemma-separable-smooth}{{00TV}{source}{}{sourcepointer-6}{}}
\newlabel{algebra-lemma-separable-smooth@cref}{{[originallemma][0][]00TV}{[1][0][]source}{}{}{}}
\makeatother
\title{Contracting an exceptional curve on an at-worst-nodal scheme keeps it at-worst-nodal of relative dimension one}
\author{The Stacks Project Authors}
\date{stacks-project:models.tex:local-snapshot-20261012; GNU FDL 1.2 or later}
\begin{document}
\maketitle
\noindent\textit{Editorial note.} The counted unit is the Stacks Project Lemma lemma-blowdown-at-worst-nodal at source lines 6311-6319 together with its proof at source lines 6321-6355, a straight adjacent pairing in the chapter file. The chapter is the upstream file models.tex; the pool directory carries the longer name stacks-models and is not used as the chapter name anywhere in this card.\par\medskip
\noindent\textit{Editorial note.} This card has no context span: the statement and proof are self-contained apart from the six pointers, and the notions they use are recorded in the standard prerequisites.\par\medskip
\noindent\textit{Editorial note.} The retained proof quotes six results of the Resolution of Surfaces, Curves, Morphisms and Algebra chapters; all six are declared as non-load-bearing navigation pointers with their stable Stacks tags looked up in the shared tag index, and none of their statements is reproduced inside the span.\par\medskip
\noindent\textit{Editorial note.} The companion card of this lane on lemma-blowup-at-worst-nodal handles the resolution by blowing up and is not touched here.\par\medskip
\noindent\textit{Editorial note.} The retained proof span starts at line 6321, the line of the opening marker read from the source, and contains exactly one proof environment.\par\medskip
\noindent\textit{Editorial note.} No citation occurs in the retained spans, so no bibliography entry is supplied.\par\medskip
\noindent\textit{Allowed standard prerequisites (editorial).}\begin{itemize}
\item At-worst-nodal schemes of relative dimension one and contractions of exceptional curves. A scheme of relative dimension one over a discrete valuation ring is at-worst-nodal when its singularities are ordinary double points whose branches are smooth, equivalently when it is smooth except at finitely many points where the local equation is the standard node; a curve over a field is nodal when it is smooth except for such double points, and smoothness of a morphism of relative dimension one at a point is tested by the regularity of the local ring combined with separability of the residue extension. An exceptional curve of the first kind is a smooth rational curve of self-intersection minus one, and it may be contracted to a regular surface. The author's own lemmas of the chapter and the quoted results of the Algebra, Curves, Morphisms and Resolution of Surfaces chapters are the allowed inputs.. The counted lemma is the blowdown half of the resolution procedure of the semistable reduction section: it keeps the class of at-worst-nodal schemes stable under the contractions performed in the resolution of the special fibre.
\end{itemize}\medskip
\ifdefined\hypertarget\hypertarget{sourcepointer-1}{}\fi
\noindent\textit{Source reference.} Original section 0C2I: focused-seeds/AGP/sources/stacks-snapshot-20261011/tags.txt, source lines 15158--15158. Allowed input of the Resolution of Surfaces chapter quoted by the retained proof as the place where exceptional curves of the first kind are defined; the pointer is registered so the reference resolves and the cited material is not reproduced inside the span. The number is the stable Stacks tag of the target, looked up in the shared tag index.\par\medskip
\ifdefined\hypertarget\hypertarget{sourcepointer-2}{}\fi
\noindent\textit{Source reference.} Original lemma 0C4D: focused-seeds/AGP/sources/stacks-snapshot-20261011/tags.txt, source lines 15223--15223. Allowed input of the Curves chapter quoted by the retained proof for the local description of a nodal curve; the pointer is registered so the reference resolves and the cited statement is not reproduced inside the span. The number is the stable Stacks tag of the target, looked up in the shared tag index.\par\medskip
\ifdefined\hypertarget\hypertarget{sourcepointer-3}{}\fi
\noindent\textit{Source reference.} Original lemma 0CD7: focused-seeds/AGP/sources/stacks-snapshot-20261011/tags.txt, source lines 15532--15532. Allowed input of the Curves chapter quoted by the retained proof for the criterion of smoothness of a morphism of relative dimension one; the pointer is registered so the reference resolves and the cited statement is not reproduced inside the span. The number is the stable Stacks tag of the target, looked up in the shared tag index.\par\medskip
\ifdefined\hypertarget\hypertarget{sourcepointer-4}{}\fi
\noindent\textit{Source reference.} Original lemma 01V9: focused-seeds/AGP/sources/stacks-snapshot-20261011/tags.txt, source lines 2438--2438. Allowed input of the Morphisms chapter quoted by the retained proof for the smoothness at a point of the total space away from the singular points; the pointer is registered so the reference resolves and the cited statement is not reproduced inside the span. The number is the stable Stacks tag of the target, looked up in the shared tag index.\par\medskip
\ifdefined\hypertarget\hypertarget{sourcepointer-5}{}\fi
\noindent\textit{Source reference.} Original lemma 00NQ: focused-seeds/AGP/sources/stacks-snapshot-20261011/tags.txt, source lines 868--868. Allowed input of the Algebra chapter quoted by the retained proof for the Cohen-Macaulay property of a regular local ring, used in the local computation; the pointer is registered so the reference resolves and the cited statement is not reproduced inside the span. The number is the stable Stacks tag of the target, looked up in the shared tag index.\par\medskip
\ifdefined\hypertarget\hypertarget{sourcepointer-6}{}\fi
\noindent\textit{Source reference.} Original lemma 00TV: focused-seeds/AGP/sources/stacks-snapshot-20261011/tags.txt, source lines 1089--1089. Allowed input of the Algebra chapter quoted by the retained proof for the passage from a regular local ring with separable residue extension to smoothness at the point; the pointer is registered so the reference resolves and the cited statement is not reproduced inside the span. The number is the stable Stacks tag of the target, looked up in the shared tag index.\par\medskip
\begin{lemma}
\label{lemma-blowdown-at-worst-nodal}
Let $R$ be a discrete valuation ring with fraction field $K$
and residue field $k$. Assume $X \to \Spec(R)$ is
at-worst-nodal of relative dimension $1$ over $R$.
Let $X \to X'$ be the contraction of an
exceptional curve $E \subset X$ of the first kind.
Then $X'$ is at-worst-nodal of relative dimension $1$ over $R$.
\end{lemma}

\begin{proof}
Namely, let $x' \in X'$ be the image of $E$.
Then the only issue is to see that $X' \to \Spec(R)$
is at-worst-nodal of relative dimension $1$
in a neighbourhood of $x'$.
The closed fibre of $X \to \Spec(R)$ is reduced, hence
$\pi \in R$ vanishes to order $1$ on $E$.
This immediately implies that
$\pi$ viewed as an element of
$\mathfrak m_{x'} \subset \mathcal{O}_{X', x'}$ but
is not in $\mathfrak m_{x'}^2$.
Since $\mathcal{O}_{X', x'}$ is regular of
dimension $2$ (by definition of contractions in
Resolution of Surfaces, Section \ref{resolve-section-minus-one}),
this implies that $\mathcal{O}_{X'_k, x'}$
is regular of dimension $1$
(Algebra, Lemma \ref{algebra-lemma-regular-ring-CM}).
On the other hand, the curve $E$ has to meet at
least one other component, say $C$ of the closed fibre $X_k$.
Say $x \in E \cap C$. Then $x$ is a node of the special
fibre $X_k$ and hence $\kappa(x)/k$ is finite separable,
see Algebraic Curves, Lemma \ref{curves-lemma-nodal}.
Since $x \mapsto x'$ we conclude that $\kappa(x')/k$
is finite separable.
By Algebra, Lemma \ref{algebra-lemma-separable-smooth}
we conclude that $X'_k \to \Spec(k)$ is smooth
in an open neighbourhood of $x'$.
Combined with flatness, this proves that
$X' \to \Spec(R)$ is smooth in a neighbourhood of $x'$
(Morphisms, Lemma \ref{morphisms-lemma-smooth-at-point}).
This finishes the proof as a smooth morphism of
relative dimension $1$ is at-worst-nodal of relative
dimension $1$
(Algebraic Curves, Lemma \ref{curves-lemma-smooth-relative-dimension-1}).
\end{proof}

\end{document}
